%%%%%%%%%%  scan idx 74  |  folio 61  %%%%%%%%%%
\SGAPrintPage{61}%
$\Atil$ subject to the restrictions $\varepsilon(\xi_i)=0$ and $\eta_j=[n_j,1+\eta_j^{(1)}+\cdots+\eta_j^{(n_j)}]$} (thus $\eta_j^{(k)}=0$ for $k>n_j$), \emph{we have}:
\[
 R((\lambda^k(\xi_i)),(\lambda^k(\eta_j)))\in\Atil^p
\]
\emph{Now let $X$ be an irreducible nonsingular quasi-projective variety. For arbitrary elements $x'_i$ of $\underline K(X)^1$ and vector bundles $E'_j$ of dimension $n_j$ on $X$ ($1\leq i\leq q$, $1\leq j\leq r$), we have}
\sgaStableTarget{sga6.e00.eq.unnamed}\[
 R((\lambda^k(x'_i)),(\lambda^k(\gamma(E'_j))))\in\underline K(X)^p\qquad\text{\SGAAuthorFootMark{*}}\ .
\]
\SGAAuthorFootText{*}{We readily prove that the hypothesis made implies that, if $R$ is expressed in terms of the $\gamma^k(\xi_i)$, $\gamma^k(E_j-\varepsilon(E_j))$, then all terms of $R$ have weight $\geq p$ (where $\gamma^k$ has weight $k$). Then 2.14 is a particular case of X 1.3.2.}

\emph{Proof}: By \S 2, the $x'_i$ themselves are differences of elements of the form $\gamma(E)$, where $E$ is a vector bundle. On the other hand, assuming $X$ locally closed in a projective space $P$, it follows from \S 2 and proposition 2.5 that, for every $F\in\cF(X)$, the sheaf $F(n)$ is generated by its sections for sufficiently large $n$. If $F$ is defined by a vector bundle $E$, this proves that $E(n)$ is the inverse image of the canonical vector bundle on a Grassmannian variety $G$ under a morphism from $X$ to $G$. Considering the corresponding morphism from $X$ to $P\times G$, we see that $E$ is the inverse image of a vector bundle on $P\times G$. Thus, applying this to all the vector bundles involved, and using corollary 2 of theorem 2.12, we see that \emph{we are reduced to the case where} $X$ \emph{is a product of Grassmannians}. In this case, the conclusion will follow from the following lemma:

\sgaStableTarget{sga6.e00.lem.2-15}\textbf{Lemma 2.15.}
\emph{Assume that $X$ is a product of Grassmannians. Then we can find a $\lambda$-homomorphism from $\underline K(X)$ to a $\lambda$-ring without torsion of the form
%%%%%%%%%%  scan idx 75  |  folio 62  %%%%%%%%%%
\SGAPrintPage{62}%
$\Atil$ (where $A$ is a graded ring), compatible with the filtrations, and such that the corresponding homomorphism $G(\underline K(X))\longrightarrow G(A)$ is injective.}

To prove the lemma, we shall use Chow's theory of Chern classes; a direct proof of the lemma, or of theorem 2.14, would yield a purely ``sheaf-theoretic'' theory of Chern classes (cf.\ \S 5).

We shall consider the homomorphism $\Ctil$ from $\underline K(X)$ to $\Atil(X)$ and put $A=A(X)$; we must prove that the homomorphism $\psi=G(\Ctil)$ from $G(\underline K(X))$ to $G(A(X))\simeq A(X)$ is injective. The fact that $\Ctil$ is compatible with the filtrations holds for every nonsingular quasi-projective variety and follows readily from the corollary to proposition 2.3. On the other hand, as $X$ is a product of Grassmannians $C_j$, we have on $X$ vector bundles $E_j$, inverse images of the canonical bundles on the factors. We shall use the following (known) facts:
\begin{enumerate}[label=\alph*)]
\item On $X$, linear equivalence of cycles is equivalent to numerical equivalence, and $A(X)$ is a group of finite rank;
\item The $C^i(E_j)$ generate the ring $A(X)$.
\end{enumerate}

From a) it follows that $\varphi:A(X)\longrightarrow G(\underline K(X))$ is injective (hence bijective), by the following lemma:

\sgaStableTarget{sga6.e00.lem.2-15.occ2}\textbf{Lemma 2.15.} \emph{Let $X^n$ be a nonsingular projective variety, and $Z^{n-p}$ a cycle on $X$ such that $\varphi(Z)\in G^p(\underline K(X))$ is zero; then $Z^{n-p}$ is numerically equivalent to zero.}

To prove this lemma, it suffices to use the additive function $\chi(F)=\sum_{i\geq0}(-1)^i\dim H^i(X,F)$ and the fact that this function equals
%%%%%%%%%%  scan idx 76  |  folio 63  %%%%%%%%%%
\SGAPrintPage{63}%
$1$
for $F=\underline{O}_{(x)}$ for every $x\in X$.

Thus, in the case considered, $A(X)$ and $G(\underline K(X))$ have the same finite rank in every degree, and $G(\underline K(X))$ has no torsion, since this holds for $A(X)$. From assertion b) above it follows, in a ``formal'' way, that the map $\Ctil$ from $\underline K(X)$ to $\widetilde{A(X)}$ is surjective, modulo finite groups. It readily follows, since $G(\underline K(X))$ and $A(X)$ have the same finite rank in every degree, that the homomorphism $G(\Ctil)$ from $G(\underline K(X))$ to $G(\widetilde{A(X)})$ is bijective modulo finite groups in every degree (use induction on the degree). As $G(\underline K(X))$ has no torsion, this homomorphism $G(\Ctil)$ is therefore injective.
\begin{flushright}q.e.d.\end{flushright}

N.B. The homomorphism $\Ctil$ is not surjective in general, even when $X$ is a projective space of dimension $2$.

\sgaStableTarget{sga6.e00.sec.5.occ2}\subsection*{\S 5. Sheaf-theoretic definition of Chern classes. Application to the study of injection morphisms.}

\sgaStableTarget{sga6.e00.thm.2-16}\textbf{Theorem 2.16.} \emph{Let $X$ be a nonsingular quasi-projective variety. Then, for every $x\in\underline K(X)$ and every integer $i\geq0$, the element $\gamma^i(x-\varepsilon(x))$ has filtration $\geq i$. Let $c^i(x)\in G^i(\underline K(X))$ be the class of this element modulo $\underline K(X)^{i+1}$; if we put:}
\[
 \widetilde c(x)=\left[\varepsilon(x),1+\sum_{i\geq1}c^i(x)\right]\in\widetilde{G(\underline K(X))},
\]
\emph{then $\widetilde c$ satisfies the conditions analogous to those of theorem 2.2.}

The first assertion follows from theorem 2.14, as does the fact that $x\longrightarrow\widetilde c(x)$ is a homomorphism of $\lambda$-rings (formula (1.25) implies $\gamma_t(x+y)=\gamma_t(x)\gamma_t(y)$, which already shows that $\widetilde c$ is an additive homomorphism). The ``functoriality'' of $\widetilde c$ is verified immediately, and
%%%%%%%%%%  scan idx 77  |  folio 64  %%%%%%%%%%
\SGAPrintPage{64}%
(2.8) follows from (2.31).
\begin{flushright}q.e.d.\end{flushright}

\sgaStableTarget{sga6.e00.cor.1.occ3}\textbf{Corollary 1:} \emph{We have $c^i(x)=\varphi(C^i(x))$.}

This follows from a uniqueness theorem for Chern classes, proved as in Hirzebruch by passage to flag varieties. We need only show that, if $Y$ is a flag variety associated with a vector bundle on $X$, then $G(\underline K(X))\longrightarrow G(\underline K(Y))$ is injective, which is verified by explicitly determining $\underline K(Y)$ with all its structures from $\underline K(X)$ (I shall not give the details here).

\sgaStableTarget{sga6.e00.cor.2.occ4}\textbf{Corollary 2:} \emph{Assume that, for every element $Z\in A(X)$, there exists a morphism $f$ from $X$ to a product $X'$ of Grassmannians such that $Z$ is in the image of $f^*$. Then the canonical homomorphisms}
\[
 \varphi:A(X)\longrightarrow G(\underline K(X)),\qquad \psi:G(\underline K(X))\longrightarrow A(X),
\]
\emph{satisfy in dimension $i$ the following relations:}
\sgaStableTarget{sga6.e00.eq.2-35}\begin{equation}\tag{2.35}
 \psi\varphi=(-1)^{i-1}(i-1)!\,\underline 1,\qquad
 \varphi\psi=(-1)^{i-1}(i-1)!\,\underline 1.
\end{equation}

As $\varphi$ is surjective, it evidently suffices to prove the first formula, and for this we can assume $X$ a product of Grassmannians. As $\varphi$ is then injective, it suffices to prove the second formula (2.35), which, by corollary 1, is written
\[
 \gamma^i(x-\varepsilon(x))\equiv(-1)^{i-1}(i-1)!\,x\qquad\mathrm{mod}\ \underline K(X)^{i+1}
\]
for every $x\in\underline K(X)^i$. As $\psi$ is \emph{injective}, it suffices to prove the analogous relation for $\Ctil(x)\in\Atil(X)^i$, and this is then a consequence of (1.22 bis).

\sgaStableTarget{sga6.e00.rem.unnamed.occ4}\textbf{Remark:} Corollary 2 makes the validity of (2.35) in all cases plausible; there should be a direct proof of it. In
%%%%%%%%%%  scan idx 78  |  folio 65  %%%%%%%%%%
\SGAPrintPage{65}%
any case, we readily prove the second formula (2.35) modulo the torsion subgroup of $G(\underline K(X))$, and, in characteristic $0$, this same formula is a consequence of the Riemann--Roch theorem. The interest of (2.35) (when it holds) is that it shows that $\varphi$ is injective modulo torsion groups (perhaps always zero!\SGAAuthorFootnote{*}{No, cf.\ the note on page 30.}), so that we ``do not lose much'' by defining the Chern classes as elements of $G(\underline K(X))$; after taking the tensor product with $\Q$, as we shall do later for the Riemann--Roch theorem, the two definitions of Chern classes would then be identical.

Let $X$ still be nonsingular and quasi-projective; let also $Y$ be a nonsingular subvariety of $X$, and $i$ the injection of $Y$ into $X$. For $i^!$ and $i_!$, we then have the following formulas:
\begin{enumerate}[label=(\roman*)]
\item $i^!$ is a homomorphism of $\lambda$-rings
\item $i_!(x.i^!(y))=i_!(x).y$, whence in particular $i_!i^!(y)=\xi y$ with $\xi=i_!i^!(1)=\gamma_X(Y)$,
\item $i^!(i_!(x))=x.\lambda_{-1}(N)$, whence in particular $i^!(\xi)=i^!i_!(1)=\lambda_{-1}(N)$,
\item $i_!(x)i_!(x')=i_!(xx'\lambda_{-1}(N))$,
\end{enumerate}
where, by definition, $N$ is the class in $\underline K(Y)$ of the conormal bundle of $Y$ in $X$. Formulas (i) and (ii) hold for arbitrary morphisms, not necessarily injections, and (iii) and (iv) are verified from definition (2.15) and the analogous formula for inverse images (taking for $x$ and $x'$ classes of vector bundles, and using the fact that
\sgaStableTarget{sga6.e00.eq.2-36}\begin{equation}\tag{2.36}
 \Tor^{\underline{O}_X}_i(\underline{O}_Y,\underline{O}_Y)=\underline{O}_Y(\wedge^i N).
\end{equation}
)

